% M05 第 4 节中文独立草案（2026-08-28）
% 用途：供逐段审阅；尚未并入 ebb2 当前中英文主稿。
% 依赖：当前 M05NEW 主稿的宏、参考文献、图件和方法部分标签。
% 权威范围：仅质量模型 A 与质量模型 B 的当前通量闭合结果。

\section{结果}
\label{sec:results}

本节沿图~\ref{fig:workflow} 的端到端流程组织结果：先建立质量模型 A 的本底与响应基线，再由其选后事例追溯本底来源；随后给出质量模型 B 的完整几何、全链选择结果和来源重分配，最后在共同时间轴上比较两套完整设计的任务性能。聚焦信号、瞬发大气本底和延迟活化本底三条物理分支分别形成探测器侧事例目录。对每套质量模型，各本底分量先按统一曝光归一生成泊松到达时刻并合并到同一有序时间轴；相邻间隔不超过 $1\,\mu\mathrm{s}$ 的到达按传递关系组成同一事例组，组内原始 TES、塑料闪烁体和 BGO 沉积分别求和，随后才依次施加探测器响应、主动 veto 和 Compton 拓扑/孔径判据。质量模型 A 对 BGO 与塑料闪烁体分别求和并各施加 50\,keV 门限；质量模型 B 不含塑料通道，只对 BGO 施加同一门限。对质量模型 A，Plastic-only 与 BGO-only 均从同一 pre-veto 母样本独立计算，只作并列的反事实诊断；combined 才是在同一时间组内同时要求两种主动通道的总沉积均低于各自门限。质量模型 B 的 combined 则与 BGO-only 退化为同一 BGO 判据。五个公共时间锚点和 81 节点任务折叠采用相同的构造与插值规则。因此，能量响应、科学窗、共轴计时与 Compton 判据保持一致；主动屏蔽通道构成则属于各自完整设计的一部分。

\subsection{质量模型 A 的基线本底与选择性能}
\label{subsec:model_a_baseline_zh}
\label{subsec:reference_background_results_zh}

质量模型 A 构成本文的基线。表~\ref{tab:phase2_cutflow} 给出第 15 天直接、无符合样本在 $[510.58,511.42)\keV$ 科学窗内的完整 cut-flow。BGO 与塑料闪烁体的联合主动门将总本底由 3086 条记录、$5.43630\cps$ 降至 443 条记录、$6.07099\times10^{-2}\cps$，相当于去除预 veto 率的 $98.883\%$。七类瞬发粒子由 $2.91217\cps$ 降为零；宽带 $\gamma$ 连续本底由 $2.35667\cps$ 降至 $1.78923\times10^{-2}\cps$；延迟活化由 $0.167460\cps$ 降至 $4.28176\times10^{-2}\cps$。

最终 Compton 拓扑/孔径选择后，本底降至 400 条记录、$5.30571\times10^{-2}\cps$，即保留主动门后带权率的 $87.394\%$。同一 $37{,}194$ 光子焦面样本中，测量窗内的 28,612 条孤立聚焦信号记录全部通过主动门；最终保留 28,006 条，对应有效面积 $15.12324\pm0.04492\,\mathrm{cm^2}$，为测量窗信号的 $97.882\%$。因此，在当前宽松策略下，主动反符合承担主要的率压低，拓扑检验移除较小的残余，同时保留绝大多数聚焦信号。

\begin{table}[H]
\centering
\caption{质量模型 A 科学窗 cut-flow。本底为第 15 天直接、无符合记录数/率；信号为共同 $37{,}194$ 光子焦面样本的记录数/有效面积。}
\label{tab:phase2_cutflow}
\footnotesize
\begin{tabular}{lccc}
\toprule
计算流 & 预 veto & 主动 veto 后 & 最终拓扑/孔径后 \\
\midrule
七类瞬发粒子 & 869 / $2.91217\cps$ & 0 / 0 & 0 / 0 \\
宽带 $\gamma$ 连续本底 & 922 / $2.35667\cps$ & 7 / $1.78923\times10^{-2}\cps$ & 6 / $1.53363\times10^{-2}\cps$ \\
延迟活化 & 1295 / $0.167460\cps$ & 436 / $4.28176\times10^{-2}\cps$ & 394 / $3.77209\times10^{-2}\cps$ \\
总本底 & 3086 / $5.43630\cps$ & 443 / $6.07099\times10^{-2}\cps$ & 400 / $5.30571\times10^{-2}\cps$ \\
聚焦信号 & 28612 / $15.45048\,\mathrm{cm^2}$ & 28612 / $15.45048\,\mathrm{cm^2}$ & 28006 / $15.12324\,\mathrm{cm^2}$ \\
\bottomrule
\end{tabular}
\end{table}

最终直接本底的非零中心值由延迟活化和宽带 $\gamma$ 连续本底构成（表~\ref{tab:reference_background_budget_zh}）。延迟活化为 $3.77209\times10^{-2}\cps$，占 $71.095\%$；宽带 $\gamma$ 连续本底为 $1.53363\times10^{-2}\cps$，占 $28.905\%$；七类瞬发粒子的中心值为零，其有限样本上限在表~\ref{tab:prompt_family_intervals_zh} 中单独报告。各分量的单事例权重不同，因此记录数不能代替有效样本量；总本底 400 条记录对应 $N_{\mathrm{eff}}=46.77$。

\begin{table}[H]
\centering
\caption{质量模型 A 第 15 天科学窗最终直接本底。标准误差和有效样本量由不等权事例的 $\sum_i w_i^2$ 计算。}
\label{tab:reference_background_budget_zh}
\footnotesize
\begin{tabular}{lrrrrr}
\toprule
分量 & 记录数 & 率 / $\cps$ & 标准误差 / $\cps$ & $N_{\mathrm{eff}}$ & 总本底 / \% \\
\midrule
延迟活化 & 394 & $3.77209\times10^{-2}$ & $4.58152\times10^{-3}$ & 67.79 & 71.095 \\
宽带 $\gamma$ 连续本底 & 6 & $1.53363\times10^{-2}$ & $6.26100\times10^{-3}$ & 6.00 & 28.905 \\
\midrule
总计 & 400 & $5.30571\times10^{-2}$ & $7.75825\times10^{-3}$ & 46.77 & 100.000 \\
\bottomrule
\end{tabular}
\end{table}

第~\ref{subsec:final_mission_performance_zh}~节从 pre-veto 物理分量重新生成泊松到达并完成共轴全链；该任务折叠给出质量模型 A 在参考线流量 $F_0=10^{-4}\phcms$ 下的 20\,d 累计信号 1637.01、本底 87632.34，高斯显著度为 5.530，20\,d 高斯 $3\sigma$ 最小可分辨流量为 $(5.425\pm0.403)\times10^{-5}\phcms$。这里的误差只是有限模板、五锚点重放和信号统计的局部传播 $1\sigma$，不等同于完整联合区间。这些任务量在该节与质量模型 B 同口径展开比较。

\subsection{质量模型 A 的本底来源与空间耦合}
\label{subsec:mass_model_a_origins_zh}

总率回答“有多少本底”，却不能回答“本底为什么入选”。为此，质量模型 A 的选后目录依次按四个层级追溯：触发事例的初级入射粒子族、产生放射性母核的材料、具体的“源体积--母核素”组合，以及该源体积相对 TES 的空间区域。这个层级把基线率、来源诊断与设计响应分开；下述绝对率和来源比例全部取自当前质量模型 A 的通量闭合目录。这样的组织顺序与 DIXE 从详细 TES 载荷出发、再追问入射粒子、产生区域和物理过程如何构成探测谱的本底研究路径一致~\cite{Tian2026DIXE}；这里借鉴的是诊断逻辑，而不是作数值类比。

按初级入射粒子族分解，质子诱发活化是最大的单项来源，最终率为 $2.48463\times10^{-2}\cps$（35 条记录），占质量模型 A 总直接本底的 $46.829\%$。其余延迟分量依次为 $\alpha$ 诱发的 $7.49239\times10^{-3}\cps$（26 条，$14.121\%$）、中子诱发的 $3.60054\times10^{-3}\cps$（11 条，$6.786\%$）、$\gamma$ 诱发的 $1.68617\times10^{-3}\cps$（260 条，$3.178\%$）、$e^+$ 诱发的 $7.02083\times10^{-5}\cps$（24 条，$0.132\%$）和 $e^-$ 诱发的 $2.52755\times10^{-5}\cps$（38 条，$0.048\%$）。这些分量之和为质量模型 A 的最终延迟率 $3.77209\times10^{-2}\cps$。

按材料分类，铜贡献选后延迟来源目录的 $70.87\%$。进一步追溯到源体积--母核素组合，主要具名来源是 L2 铜支撑环中的 $^{62}$Cu、50\,mK 混合室铜板中的 $^{62}$Cu 和 L0 TES 热沉环中的 $^{62}$Cu；其后的具名组包括 Bi 上半圆柱中的 $^{199}$Po，以及 50\,mK 混合室铜板和 L2 铜支撑环中的 $^{64}$Cu。这里的排序已经包含输运、响应和全部选后条件，不能由材料总质量或总活度直接推出。

空间分解进一步说明了几何耦合：DR 混合室及多级冷盘贡献选后延迟目录的 $32.18\%$，TES 近场结构贡献 $47.29\%$，两者合计 $79.47\%$。按通量闭合的最终延迟率换算，前者约为 $1.21398\times10^{-2}\cps$，后者约为 $1.78397\times10^{-2}\cps$。因此，质量模型 A 的主要延迟负担不是“铜很多”这一全局事实本身，而是部分活化铜及近场结构仍与 TES 入选条件保持较强的空间耦合。

这一来源诊断提出了一个明确但尚非因果证明的设计问题：若改变焦面与冷端结构的相对位置，并用局部主动屏蔽重新包围探测器，能否在保留聚焦信号响应的同时削弱这部分选后耦合？质量模型 B 是对该完整设计问题的独立实现。由于它同时改变物理入口、局部 W 结构、BGO 布局和部分材料清单，后文只能评价这套完整构型的条件性能，不能把任一差值单独归因于侧烟囱或焦面位移。

\FloatBarrier

\subsection{由来源诊断到质量模型 B 的完整构型}
\label{subsec:optimized_shield_results_zh}

质量模型 B 将六层 TES 沿局部 $z'=-2.80$\,cm 的侧烟囱轴布置。烟囱 Al 壁的内、外半径分别为 10.75 和 11.05\,cm，壁厚为 3\,mm，并与稀释制冷机外壳形成连续接口；烟囱底部与稀释制冷机底部之间保留 0.25\,cm 间隙。焦面由外、内方孔分别为 6.00 和 5.40\,cm、轴向深度为 2.00\,cm 的四边 W 框支撑。三个主动 BGO 体积覆盖烟囱侧面、前端光学环和后端冷孔环；质量模型 B 不含塑料闪烁体通道。主动门仍按同一 $1\,\mu\mathrm{s}$ 公共时间组累计全部 BGO 沉积，并要求组内 BGO 总沉积严格低于 50\,keV。

图~\ref{fig:optimized_shield_background_zh} 以相同仪器坐标和比例比较两套冷端结构，并在局部放大图中显示 TES、关键铜支撑、W 孔径/框架和主动 BGO 的相对位置。最终选择后，质量模型 A 与质量模型 B 的聚焦信号有效面积分别为 $15.12324\pm0.04492$ 和 $15.08544\pm0.04503\,\mathrm{cm^2}$，即质量模型 B 对当前聚焦焦面样本保留了质量模型 A 有效面积的 $99.75\%$。

\begin{figure}[!ht]
\centering
\includegraphics[width=\linewidth]{figures/section4_current/fig09_geometry_response_detailed.pdf}
\caption{质量模型 A 与质量模型 B 的冷端--焦面 $x'$--$z'$ 切面及选后延迟源位置。（a）、（b）采用完全相同的坐标范围和比例，显示多级铜冷盘、六层 TES、主动 BGO、W 孔径/框架以及质量模型 B 的侧烟囱和铜冷指；（c）、（d）以相同物理尺度放大两个焦面邻域，并在质量模型 A 中标出 L2 铜支撑环、L0 铜热沉环和 50\,mK MXC 铜盘。点为第 15 天最终入选延迟事例的放射性母核产生位置，颜色表示源材料，点面积在两模型间采用共同标度并与当前带权率成正比。质量模型 B 的来源点按粒子族分别将一阶、二阶权重矩重标到当前正式 cut-flow；该图是与输运几何尺寸对齐的语义切面，小型 CSG 开孔和服务件未逐一绘出。}
\label{fig:optimized_shield_background_zh}
\end{figure}

\begin{table}[!ht]
\centering
\caption{质量模型 B 侧烟囱在局部仪器坐标中的几何参数。表中“作用”表示设计意图，不是单因素因果测量。}
\label{tab:final_geometry_zh}
\footnotesize
\begin{tabular}{p{0.23\linewidth}p{0.34\linewidth}p{0.34\linewidth}}
\toprule
部件 & 最终设置 & 设计作用 \\
\midrule
Al 烟囱 & 内/外半径 10.75/11.05 cm；壁厚 3 mm & 将 TES 横向布置于多级冷端区域之外 \\
共同轴与间隙 & $z'=-2.80$ cm；烟囱/DR 底部 $-13.85/-14.10$ cm & 连续接口并保留 0.25 cm 间隙 \\
TES 焦面 & 沿烟囱轴布置六层 Ta 像素 & 保持探测器响应和多层阻止深度 \\
W 焦面框 & 外/内方孔 6.00/5.40 cm；深 2.00 cm & 紧凑局部支撑及孔径定义 \\
冷端/光学孔 & 1.50 cm 冷孔；6.00 cm 上端光学开孔 & 保留低温服务和聚焦光束通道 \\
主动 BGO & 侧屏、前端光学环、后端冷孔环 & 环绕烟囱并按共时间组累计的反符合 \\
BGO 阈值 & 50 keV 离线沉积能 & 共同主动 veto 判据 \\
\bottomrule
\end{tabular}
\end{table}

有效面积的一致只约束当前聚焦光子样本，并不意味着两套几何具有相同的弥散光子接受度。质量模型 A 使用半径约 1.898\,cm 的入口包络和 624-copy W 多孔准直结构；质量模型 B 使用半径 2.70\,cm 的物理窗口和开放四边 W 框，主动屏蔽构型也不同。因此，本节比较的是共同源、响应、计时和选择合同下的两套完整设计。质量模型 B 的本底和灵敏度必须以其当前入口、W 框、BGO 与无塑料通道的组合为条件；现有数据没有提供仅改变侧烟囱、同时冻结其余结构的消融样本。

\FloatBarrier

\subsection{质量模型 B 的全链本底与来源重分配}
\label{subsec:optimized_fullchain_zh}

质量模型 B 的探测器侧事例目录沿用上述 420\,eV FWHM 的 TES 响应、$[510.58,511.42)\keV$ 科学窗、主动 BGO 门以及同一 Compton 拓扑/孔径检验。表~\ref{tab:optimized_cutflow_zh} 给出第 15 天直接、无符合样本的完整 cut-flow。

主动 BGO 将总本底由 7684 条记录、$2.38372\cps$ 降至 137 条记录、$1.11379\times10^{-2}\cps$，即只剩预 veto 率的 $0.4672\%$。测量窗内的 28,434 条孤立聚焦信号记录全部通过主动门。随后，Compton 拓扑/孔径检验将本底由 137 条降至 126 条，保留主动门后带权率的 $96.492\%$；聚焦信号由 28,434 条降至 27,936 条，对应最终有效面积 $15.08544\pm0.04503\,\mathrm{cm^2}$。因此，与质量模型 A 相同，主动反符合承担主要的率压低，最终拓扑检验只移除较小的残余，同时保留 $98.249\%$ 的测量窗聚焦信号。

\begin{table}[H]
\centering
\caption{质量模型 B 经响应卷积的科学窗 cut-flow。本底为第 15 天直接、无符合记录数/率；信号为记录数/有效面积。}
\label{tab:optimized_cutflow_zh}
\footnotesize
\begin{tabular}{lccc}
\toprule
计算流 & 预 veto & 主动 BGO 后 & 最终拓扑/孔径后 \\
\midrule
七类瞬发粒子 & 1448 / $0.878692\cps$ & 3 / $1.81131\times10^{-3}\cps$ & 3 / $1.81131\times10^{-3}\cps$ \\
宽带 $\gamma$ 连续本底 & 2987 / $1.37944\cps$ & 12 / $5.35535\times10^{-3}\cps$ & 12 / $5.35535\times10^{-3}\cps$ \\
延迟活化 & 3249 / $0.125586\cps$ & 122 / $3.97120\times10^{-3}\cps$ & 111 / $3.58054\times10^{-3}\cps$ \\
总本底 & 7684 / $2.38372\cps$ & 137 / $1.11379\times10^{-2}\cps$ & 126 / $1.07472\times10^{-2}\cps$ \\
聚焦信号 & 28434 / $15.35436\,\mathrm{cm^2}$ & 28434 / $15.35436\,\mathrm{cm^2}$ & 27936 / $15.08544\,\mathrm{cm^2}$ \\
\bottomrule
\end{tabular}
\end{table}

质量模型 B 的最终直接本底由七类瞬发粒子 $1.81131\times10^{-3}\cps$、宽带 $\gamma$ 连续本底 $5.35535\times10^{-3}\cps$ 和延迟活化 $3.58054\times10^{-3}\cps$ 构成，分别占总率的 $16.854\%$、$49.830\%$ 和 $33.316\%$；总直接率为 $1.07472\times10^{-2}\cps$。与质量模型 A 相比，延迟活化和宽带 $\gamma$ 分别降至其 $9.492\%$ 和 $34.920\%$，总率降至 $20.256\%$，即降低 4.94 倍。质量模型 A 的瞬发七族中心值为零，因而不能用中心值为这项构造跨模型比值。

三项本底的相对权重随完整构型改变：质量模型 A 由延迟活化主导，而质量模型 B 中宽带 $\gamma$ 连续本底成为最大分量，并出现由三个正电子族模板支撑的瞬发中心值。这是有限统计下的来源重分配，而不是对正电子本底的稳健识别。由于入口、W 框、BGO 与局部材料同时变化，当前 A/B 对照仍不能把总率或组成差异唯一归因于其中某个结构。

\begin{table}[!ht]
\centering
\caption{质量模型 B 第 15 天最终直接本底的有限输运支撑。标准误差和有效样本量由不等权事例的 $\sum_i w_i^2$ 计算。}
\label{tab:optimized_component_precision_zh}
\scriptsize
\resizebox{\linewidth}{!}{%
\begin{tabular}{lrr@{\hspace{1.4em}}rrr}
\toprule
分量 & 记录数 & 率 / $\cps$ & 标准误差 / $\cps$ & $N_{\mathrm{eff}}$ & 总本底 / \% \\
\midrule
七类瞬发粒子 & 3 & $1.81131\times10^{-3}$ & $1.04576\times10^{-3}$ & 3.00 & 16.854 \\
宽带 $\gamma$ 连续本底 & 12 & $5.35535\times10^{-3}$ & $1.55631\times10^{-3}$ & 11.84 & 49.830 \\
延迟活化 & 111 & $3.58054\times10^{-3}$ & $5.36971\times10^{-4}$ & 44.46 & 33.316 \\
\midrule
总本底 & 126 & $1.07472\times10^{-2}$ & $1.95040\times10^{-3}$ & 30.36 & 100.000 \\
\bottomrule
\end{tabular}}
\end{table}

表~\ref{tab:optimized_component_precision_zh} 给出三项分量的不等权统计支撑。最终样本共 126 条记录，对应 $N_{\mathrm{eff}}=30.36$。在 W2 科学窗最终 Compton/孔径选择下，七类瞬发粒子保留的 3 个有限输运模板均属于正电子族，其中心率为 $1.81131\times10^{-3}\cps$，Garwood 中央 95\% 区间为 $[3.735\times10^{-4},5.293\times10^{-3}]\cps$。这个区间较宽，故结果只说明当前样本出现了三个保留模板，不能据此确认稳定的正电子本底分量。

\begin{table}[!ht]
\centering
\caption{第 15 天科学窗最终选择后七类瞬发粒子的有限样本约束。括号内为保留模板数；零计数项给出单侧 95\% 上限，模型 B 的正电子项给出中心率及 Garwood 中央 95\% 区间。}
\label{tab:prompt_family_intervals_zh}
\scriptsize
\begin{tabular}{lcc}
\toprule
粒子族 & 质量模型 A / $\cps$ & 质量模型 B / $\cps$ \\
\midrule
$\alpha$ & $<1.023\times10^{-2}$ (0) & $<1.902\times10^{-3}$ (0) \\
$e^-$ & $<9.690\times10^{-3}$ (0) & $<1.817\times10^{-3}$ (0) \\
$e^+$ & $<9.689\times10^{-3}$ (0) & $1.811\times10^{-3}\,[3.735\times10^{-4},5.293\times10^{-3}]$ (3) \\
$\mu^-$ & $<9.212\times10^{-3}$ (0) & $<1.708\times10^{-3}$ (0) \\
$\mu^+$ & $<8.897\times10^{-3}$ (0) & $<1.703\times10^{-3}$ (0) \\
$n$ & $<1.303\times10^{-2}$ (0) & $<1.818\times10^{-3}$ (0) \\
$p$ & $<9.747\times10^{-3}$ (0) & $<1.823\times10^{-3}$ (0) \\
\bottomrule
\end{tabular}
\end{table}

按当前逐族一阶、二阶权重矩重标后，质量模型 B 的铜来源为 $(2.26708\pm0.42826)\times10^{-3}\cps$，占其最终延迟率的 $63.32\%$，说明铜活化并未消失。DR/MXC 及多级冷盘来源则由质量模型 A 的 $(1.21398\pm0.26869)\times10^{-2}\cps$ 降至质量模型 B 的 $(1.40604\pm1.09289)\times10^{-4}\cps$；中心值相差 86.3 倍，但质量模型 B 的这一小分量只有 7 条来源记录，因而相对统计误差较大。与此同时，总延迟率由 $(3.77209\pm0.45815)\times10^{-2}$ 降至 $(3.58054\pm0.53697)\times10^{-3}\cps$。质量模型 B 的主要具名组仍包括 L5 铜环和 TES 热沉中的 $^{62}$Cu。选后铜贡献因而发生削弱和空间重分配，而不是材料层面的完全消除。

\begin{figure}[!ht]
\centering
\includegraphics[width=\linewidth]{figures/section4_current/fig09b_geometry_response_rates.pdf}
\caption{来源诊断与完整构型响应的定量比较。（a）质量模型 A、B 的第 15 天最终延迟率按冷端、TES 近场和其他结构分解；横向误差棒为 $\sqrt{\sum_i w_i^2}$。质量模型 B 的来源事件按粒子族分别重标一阶率矩和二阶方差矩，使两者严格闭合到当前正式延迟 cut-flow。（b）质量模型 B/A 的冷端延迟、铜来源延迟、总延迟、科学窗总本底和聚焦信号有效面积比；虚线表示两模型相等。误差为有限输运样本的局部独立 $1\sigma$ 传播，不包含联合系统误差。}
\label{fig:geometry_response_rates_zh}
\end{figure}

图~\ref{fig:s33_compton_mult} 给出最终本底的 TES 命中多重度。480--550\,keV 最终样本在质量模型 A 和 B 中分别含 566 和 776 条本底记录；两者均只有 1 条五命中记录，六命中及以上均为零。五命中带权率分别为 $2.87731\times10^{-6}$ 和 $6.93521\times10^{-5}\cps$，只占相应宽窗最终率的 $0.002372\%$ 和 $0.105849\%$。科学窗内，两套设计的最高非零多重度分别为四和三，五命中及以上均为零。最终科学窗事例仍以单像素为主：质量模型 A 的单/多像素贡献为 293/107 条、$4.19886\times10^{-2}/1.10685\times10^{-2}\cps$，质量模型 B 为 93/33 条、$8.82597\times10^{-3}/1.92123\times10^{-3}\cps$。

\begin{figure}[!ht]
\centering
\includegraphics[width=\linewidth]{figures/section4_current/fig07_hit_multiplicity.pdf}
\caption{质量模型 A 与质量模型 B 第 15 天最终选择后带权本底率按 TES 命中多重度的分解：（a）480--550 keV 宽带；（b）$[510.58,511.42)\keV$ 科学窗。误差棒为各多重度分箱有限输运权重的 $\sqrt{\sum_i w_i^2}$；四组最终样本均无 $n\geq6$ 的事例。}
\label{fig:s33_compton_mult}
\end{figure}

480--550\,keV 宽窗提供独立于窄窗的选择诊断。质量模型 A 的宽窗带权率由预 veto 的 $10.7984\cps$ 降至联合主动门后的 $0.140797\cps$，最终为 $0.121285\cps$；质量模型 B 相应由 $4.92157\cps$ 降至 $0.0735321\cps$，最终为 $0.0655197\cps$。这与科学窗 cut-flow 给出相同的阶段排序：主动反符合承担主要抑制，Compton 拓扑/孔径检验对残余作较小修正。

\begin{figure}[!ht]
\centering
\includegraphics[width=\linewidth]{figures/section4_current/fig05_day15_selection_spectra_current.pdf}
\caption{第 15 天质量模型 A、B 的 TES 测量能量微分计数率谱。（a）、（b）为 480--550\,keV 宽带，当前 0.25\,keV 诊断分箱按四箱合并为 1\,keV；（c）、（d）保留 508--514\,keV 内原生 0.25\,keV 分箱，灰带及虚线标出 $[510.58,511.42)\keV$ 科学窗。三条曲线依次为主动 veto 前、联合主动 veto 后和最终 Compton 拓扑/孔径选择后；宽带面板中末级选择用半透明带表示 $\sqrt{\sum_i w_i^2}$，局部面板给出逐箱误差棒。零权重分箱保持断开，不加入平滑或伪计数。该图表示探测器计数率谱，而非入射环境通量谱。}
\label{fig:day15_selection_spectra_zh}
\end{figure}

上述结果闭合了“基线总率--来源诊断--候选构型--重新选后”的探测器侧链条：在当前完整构型和分析合同下，质量模型 B 的延迟活化与宽带 $\gamma$ 连续本底中心率均低于质量模型 A，并在扩大的瞬发目录中留下三个正电子族模板；最终组成也由质量模型 A 的延迟活化主导转为质量模型 B 的宽带 $\gamma$ 连续本底占比最高。由于入口、W 结构、主动屏蔽和局部材料同时变化，这一闭合支持质量模型 B 作为完整候选设计的条件性能评价，但不构成侧烟囱这一单一几何因素的因果测量。

\FloatBarrier

\subsection{共同时间轴与任务性能}
\label{subsec:mission_fold_validation_zh}
\label{subsec:final_mission_performance_zh}

任务投影前，逐粒子族瞬发环境折叠在三个非参考环境中做过直接输运检验。L1、H1 和 L2 的直接输运率/解析率分别为 0.992、1.004 和 0.994，最大偏差为 $0.51\sigma$。该检验只覆盖 $e^+$、中子和 $\gamma$ 三族在 480--550\,keV 内的瞬发切片；它不验证当前窄科学窗、活化产生或延迟衰变分量，因而不对这些分支作超出样本的外推。

表~\ref{tab:phase2_cutflow} 和表~\ref{tab:optimized_cutflow_zh} 的“直接、无符合”cut-flow 只作单例诊断，不是任务生产顺序。第~\ref{sec:mission_time_fold}~节定义的五个锚点均从 pre-veto 物理分量率重新生成泊松到达，在共同时间轴上依次完成传递分组、原始通道组和、响应、veto 与 Compton；同阶段的“共轴最终率/直接无符合最终率”定义本底修正 $\widehat\rho_g$，聚焦信号重放则给出偶然存活率 $\widehat\eta_g$。两者从第 0、5、10、15 和 20 天线性插值到 81 个任务节点。图~\ref{fig:common_time_anchors} 给出这些系数及其对直接加和率的修正。五个锚点上，质量模型 A 的 $\widehat\rho_g$ 依次为 0.9307、0.9499、0.9708、0.9509 和 1.0060，质量模型 B 为 1.0293、1.0026、0.9905、0.9951 和 0.9901。各设计在五个锚点分别采用固定曝光和相同的分组规则；81 节点曲线只是五点线性插值，不代表 81 次独立泊松重放。

\begin{figure}[H]
\centering
\includegraphics[width=\linewidth]{figures/section4_current/fig04_common_time_normalization.pdf}
\caption{五锚点公共时间诊断。（a）质量模型 A 与质量模型 B 的科学窗最终直接率和按既定系数修正的率；锚点误差棒为固定曝光下传播到修正率的泊松标准误差。（b）本底修正 $\rho_g(t)$ 与聚焦信号存活率 $\eta_g(t)$，后者误差棒为二项统计误差。质量模型 A、B 的每点等效曝光分别为 $2\times10^4$ 和 $2\times10^5$\,s，所有锚点采用同一 $1\,\mu\mathrm{s}$ 分组规则。曲线是五点的线性插值，因此点线闭合只检验既定系数及其插值，不是环境外推的独立验证。}
\label{fig:common_time_anchors}
\end{figure}

聚焦信号的 $\widehat\eta_g$ 在质量模型 A 中为 0.96810--0.97039，在质量模型 B 中为 0.99547--0.99580；第 15 天分别为 0.9682365 和 0.995472。质量模型 B 的较高值与其当前完整构型下较低的符合占有率一致，而不是对聚焦信号采用了不同的选择判据。固定的 $\rho_g$ 和 $\eta_g$ 插值到 0--20\,d、步长 0.25\,d 的 81 个节点后，与同一轨迹透过率和分量响应作梯形积分。

20\,d 时，质量模型 A 累计 1637.01 个参考流量信号和 87632.34 个本底计数；质量模型 B 累计 1678.39 个信号和 18568.89 个本底计数。任务产物在这一层保留“宽带 $\gamma$ 连续本底”和“非连续谱本底”两项加和：质量模型 A 分别为 25760.39 和 61871.95 个计数，质量模型 B 为 9359.44 和 9209.45 个；后者同时包含七类瞬发粒子和延迟活化，不能整体解释为活化。质量模型 B 对应的累计本底中心值比质量模型 A 少 69063.45 个计数，为后者的 $21.19\%$。

图~\ref{fig:optimized_mission_significance_zh} 和表~\ref{tab:primary_sensitivity_zh} 汇总共同分析合同下的条件任务投影。20\,d 高斯显著度在质量模型 A、B 中分别为 5.530 和 12.317，Asimov 值分别为 5.513 和 12.138。高斯 $3\sigma$ 最小可分辨流量分别为 $(5.425\pm0.403)\times10^{-5}$ 和 $(2.436\pm0.223)\times10^{-5}\phcms$；对应 Asimov 值分别为 $(5.434\pm0.403)\times10^{-5}$ 和 $(2.445\pm0.223)\times10^{-5}\phcms$。这些误差为局部传播 $1\sigma$；高斯近似与 Asimov 计算给出相同的中心值排序，但均未包含完整联合不确定度。

\begin{figure}[!ht]
\centering
\includegraphics[width=0.98\linewidth]{figures/section4_current/fig10_mission_performance.pdf}
\caption{质量模型 A 与质量模型 B 在共同分析合同下的 81 节点条件任务性能。（a）$F_0=10^{-4}\phcms$ 下的高斯计数显著度，水平线标出 $3\sigma$ 与 $5\sigma$；（b）高斯和 Asimov $3\sigma$ 最小可分辨流量。两条曲线采用相同源、探测器响应、事例选择、计时及大气透过约定，但两套完整设计的入口、局部 W 结构和主动屏蔽布局并不相同。}
\label{fig:optimized_mission_significance_zh}
\end{figure}

\begin{table}[!ht]
\centering
\caption{质量模型 A 与质量模型 B 在共同分析合同下的条件任务投影；流量阈值后的误差为局部传播的 $1\sigma$ 统计不确定度，不是联合置信区间。}
\label{tab:primary_sensitivity_zh}
\footnotesize
\resizebox{\linewidth}{!}{%
\begin{tabular}{lrr}
\toprule
量 & 质量模型 A & 质量模型 B \\
\midrule
选后有效面积 / $\mathrm{cm^2}$ & $15.12324\pm0.04492$ & $15.08544\pm0.04503$ \\
第 15 天直接本底 / $\cps$ & $5.30571\times10^{-2}$ & $1.07472\times10^{-2}$ \\
第 15 天 $F_0$ 信号 / $\cps$ & $9.54765\times10^{-4}$ & $9.79169\times10^{-4}$ \\
20 d 本底计数 & 87632.34 & 18568.89 \\
20 d $F_0$ 信号计数 & 1637.01 & 1678.39 \\
$Z_{20\mathrm d}$，高斯 / Asimov & 5.530 / 5.513 & 12.317 / 12.138 \\
$T_{3\sigma}$，高斯 / Asimov / d & 5.568 / 5.597 & 0.904 / 0.933 \\
$T_{5\sigma}$，高斯 / Asimov / d & 15.980 / 16.056 & 2.513 / 2.645 \\
$F_{3\sigma}(20\,\mathrm d)$，高斯 / Asimov / $\phcms$ &
$(5.425\pm0.403)\times10^{-5}$ / $(5.434\pm0.403)\times10^{-5}$ &
$(2.436\pm0.223)\times10^{-5}$ / $(2.445\pm0.223)\times10^{-5}$ \\
$F_{5\sigma}(20\,\mathrm d)$，高斯 / Asimov / $\phcms$ &
$(9.042\pm0.671)\times10^{-5}$ / $(9.067\pm0.671)\times10^{-5}$ &
$(4.059\pm0.372)\times10^{-5}$ / $(4.084\pm0.372)\times10^{-5}$ \\
\bottomrule
\end{tabular}}
\end{table}

\FloatBarrier
\enlargethispage{\baselineskip}

高斯阈值的变化可以直接拆成累计本底和累计信号核两项。质量模型 B 的 20\,d 本底为质量模型 A 的 $21.19\%$，而其 20\,d 累计信号核为质量模型 A 的 1.025278 倍，因此
\begin{equation}
 \frac{F_{\min,A}}{F_{\min,B}}
 =\sqrt{\frac{N_{b,A}}{N_{b,B}}}\frac{K_B^{20\mathrm d}}{K_A^{20\mathrm d}}
 =2.172397\times1.025278
 =2.227311 .
 \label{eq:fmin_ratio_decomposition_zh}
\end{equation}
这里的 $K_g^{20\mathrm d}$ 是 81 节点积分后的累计信号核，而不是某一节点的条件信号核。式~\eqref{eq:fmin_ratio_decomposition_zh} 表明，当前中心值下 2.227 倍的流量阈值差异主要来自总本底的平方根项，并由当前时间轴上的累计信号核差异进一步修正。

本轮采用 \texttt{NO\_TOPUP} 边界：不启动新的模拟，只穷尽并纳入现有已验收样本。该边界是生产与方法范围的停止决定，不由局部流量阈值相对标准误差、零条入选记录或七族 exact 区间推导。有限模板、五锚点 replay MC、信号效率以及源、响应和仪器系统学的联合传播仍未完成。

综上，质量模型 B 保留了质量模型 A 的 $99.75\%$ 聚焦有效面积，并在当前中心值下将 20\,d 累计本底和高斯 $3\sigma$ 流量阈值分别降至 $21.19\%$ 和 $44.90\%$。这是入口、W 框、主动屏蔽及局部材料同时变化的条件比较，不能识别单一几何因素，也不能解释为无条件性能优势。结果仅覆盖探测器--低温系统和合成轨迹；吊舱/光学、真实可见性、50/80\,keV 离线/本机触发门限、Laue PSF 及辐射场、截面、标定、重建、指向系统学均不在当前联合不确定度内。

\FloatBarrier
